\section{DD2：一般函数补偿}\label{sec:dd2}

本节将线性残差界推广为一般函数界，并把统一收缩比例改为指定下降函数
$\varphi(t)<t$。先给出分段补偿关系，再讨论距离下降与实际长度之间的联系。

\subsection{三段补偿定理}

沿用前节的记号。设 $\Phi:[0,R]\to\R$ 有限且 $\Phi(0)=0$，
$\psi:D_\psi\to[0,\infty)$ 有限，$0\in D_\psi\subseteq[0,\infty)$。
在所讨论的每个允许近端对上，设
\begin{equation}\label{eq:dd2-assumptions}
 \begin{gathered}
 y\in\mathcal J(x)\subseteq J_\lambda(x),\quad
 t=d(x),\quad a=d(y),\quad s=\norm{x-y},\quad \rho=r_F(y)\in D_\psi,\\
 a\le\psi(\rho),\qquad
 \exists p\in P_S(x):\quad \ip{x-y}{y-p}\ge\Phi(t).
 \end{gathered}
\end{equation}
此处不预设 $\psi$ 单调，不能将 $\psi(\rho)$ 换成 $\psi(s/\lambda)$。
指定 $\varphi:[0,R]\to[0,R]$，满足 $\varphi(0)=0$、
$0\le\varphi(t)<t$（$t>0$），并记
\begin{equation}\label{eq:dd2-normalization}
 r(t)=\frac{\varphi(t)}t,\qquad
 z(t)=\frac{\lambda}{t}\inf\{\rho\in D_\psi:\psi(\rho)>\varphi(t)\}.
\end{equation}
对 $0\le r<1$，定义
\begin{equation}\label{eq:dd2-sigma}
 \boxed{\Sigma(r,z)=
 \begin{cases}
 (r^2+z^2-1)/2,&1-r\le z\le1+r,\\
 z(z-1),&z>1+r,\\
 r(r-1),&0\le z<1-r,\quad r\ge1/2.
 \end{cases}}
\end{equation}
不在区域 $r<1/2,\ z<1-r$ 内定义 $\Sigma$。

\begin{theorem}[DD2]\label{thm:dd2}
假设输入距离在 $[0,R]$ 内的每个允许近端对满足~\eqref{eq:dd2-assumptions}。
对每个 $t>0$，若 $z(t)<\infty$，要求 $(r(t),z(t))$ 属于 $\Sigma$ 的定义域，且
\begin{equation}\label{eq:dd2-condition}
 \boxed{-\frac{\Phi(t)}{t^2}\le
 \Sigma\left(\frac{\varphi(t)}t,\
 \frac{\lambda}{t}\inf\{\rho\in D_\psi:\psi(\rho)>\varphi(t)\}\right).}
\end{equation}
仅有下述端点例外：在 $z(t)>1+r(t)$ 时，若式~\eqref{eq:dd2-normalization}
中的有限下确界属于所定义的严格超水平集，则式~\eqref{eq:dd2-condition}须用严格不等号。
若 $z(t)=\infty$，不需要该尺度的补偿条件。于是
\begin{equation}\label{eq:dd2-conclusions}
 d(y)\le\varphi(d(x)),\qquad
 \norm{x-y}\le B_\Phi(d(x)),\qquad
 B_\Phi(t)=\frac{t+\sqrt{t^2-4\Phi(t)}}2.
\end{equation}
有近端对存在的尺度上根号自动有定义；输入距离为零时 $y=x$。
\end{theorem}

\begin{proof}
由同一最近点展开，得
\begin{equation}\label{eq:dd2-basic}
 a^2+s^2\le t^2-2\Phi(t),\qquad
 s^2+\norm{y-p}^2\le t^2-2\Phi(t),\qquad \lambda\rho\le s.
\end{equation}
引理~\ref{dd1:geometry} 给出 $s,\norm{y-p}\le B_\Phi(t)$，
并由 $\Phi(0)=0$ 得零输入驻定。
固定 $t>0$，记 $c=\varphi(t)$、$\phi=\Phi(t)$、
$q=\lambda\inf\{\psi>c\}$。若 $a>c$，则 $\psi(\rho)>c$，从而 $s\ge\lambda\rho\ge q$。
$q=\infty$ 时已不可能；有限时，以下任一条件均可排除该步：
\[
 \begin{array}{ll}
 c\ge t/2,\quad \phi\ge c(t-c),& B_\Phi(t)\le c;\\
 \phi\ge(t^2-c^2-q^2)/2,& a^2+s^2>c^2+q^2;\\
 q\ge t/2,\quad \phi>q(t-q),& B_\Phi(t)<q.
 \end{array}
\]
第三行在 $\phi=q(t-q)$ 时，若 $q/\lambda\notin\{\psi>c\}$，
仍有 $s\ge\lambda\rho>q=B_\Phi(t)$。
将三个下界除以 $t^2$，写成
$A=r(1-r)$、$E=(1-r^2-z^2)/2$、$C=z(1-z)$。由
\[
 A-E=\frac{z^2-(1-r)^2}{2},\qquad
 C-E=\frac{r^2-(z-1)^2}{2}
\]
比较适用的最小下界，依次得到式~\eqref{eq:dd2-sigma}的三支。
$z=1\pm r$ 处由第二行包含等号，故唯一严格端点正如定理所述。
\end{proof}

第一支来自平方距离恒等式，第二支来自步长上界，第三支利用正的锚定下界。
$\Phi$ 保留符号，第三支的信息才能保留。
若 $D_\psi=[0,\infty)$，且 $\psi$ 连续、严格递增、无界并满足 $\psi(0)=0$，则
$z(t)=\lambda\psi^{-1}(\varphi(t))/t$。
取 $\psi(\rho)=\ell\rho$、$\Phi\le0$，即恢复 DD1 的两支公式。

固定锚定下界后，也可反求残差正则性条件。固定 $t>0$、$\phi=\Phi(t)\le t^2/4$，
记 $E=t^2-2\phi$、$B=B_\Phi(t)$。给定 $(t-B)_+\le c<B$，
若
\begin{equation}\label{eq:dd2-dual-cap}
 \psi(\rho)\le c\quad
 \text{对所有 }\rho\in D_\psi\text{ 满足 }\lambda\rho\le B,\
 \lambda\rho<\sqrt{E-c^2},
\end{equation}
则 $a\le c$。否则 $a>c$ 与式~\eqref{eq:dd2-basic}会使实际残差恰落在所列区域，
矛盾。这就由给定的锚定下界确定了残差正则性所需满足的条件。

\subsection{一般函数的扩张与长度边界}

\begin{example}[次线性收敛]\label{ex:dd2-rotation}
令 $Q(u,v)=(-v,u)$，取 $1<m<2$，定义
\[
 F(y)=\norm y^{m-1}Qy,\qquad S=\{0\},\qquad \lambda=1.
\]
每个输入 $x$ 有唯一输出：$a\ge0$ 由 $t^2=a^2+a^{2m}$ 确定，且
\[
 y=\frac{I-a^{m-1}Q}{1+a^{2m-2}}x,\qquad
 r_F(y)=s=a^m,\qquad a=r_F(y)^{1/m},\qquad \ip{x-y}{y}=0.
\]
半径方程的右侧连续严格递增且无界，故确实给出完整预解式。
取 $\Phi=0$、$\psi(\rho)=\rho^{1/m}$，以半径方程的逆函数为 $\varphi$。
此时 $r=a/t$、$z=a^m/t$，有 $r^2+z^2=1$，
故 $1-r<z<1+r$，DD2 第一支取等。

所有实际半径递减趋零。记 $\nu=2m-2$，由半径递推得
\[
 t_{k+1}^{-\nu}-t_k^{-\nu}\longrightarrow\nu/2,\qquad
 t_k\sim(\nu k/2)^{-1/\nu}.
\]
前式由 $(1+v)^{-\nu/2}=1-(\nu/2)v+o(v)$ 得到，后式由平均增量得到。
因而收敛为次线性，且 $a/r_F(y)=a^{1-m}\to\infty$，
不存在 DD1 所需的局部线性残差常数。
它的有限长度将在本节末由耦合势推出。
\end{example}

\begin{example}[低于二次阶的负缺陷]\label{ex:dd2-shear}
取 $\lambda=1$、$F(u,v)=(-\sqrt{|v|},7v)$、$S=\R\times\{0\}$。直接解得
\[
 J_1(\xi,\eta)=\left(\xi+\sqrt{|\eta|/8},\,\eta/8\right).
\]
记 $t=|\eta|$，最近输入点为 $p=(\xi,0)$，则
\[
 a=t/8,\quad \rho^2=s^2=t/8+49t^2/64,\quad
 \ip{x-y}{y-p}=-t/8+7t^2/64\ge-t/8.
\]
故 $\psi(\rho)=\rho^2$、$\Phi(t)=-t/8$ 有效。
在 $0<t\le1/8$ 上取 $\varphi(t)=t/2$，则 $r=1/2$、
$z=1/\sqrt{2t}>1+r$，且
$\Sigma(r,z)-1/(8t)=3/(8t)-1/\sqrt{2t}\ge0$。
DD2 因而允许不能由有界二次系数控制的负缺陷。
实际纵坐标以 $8^{-k}$ 衰减，横向增量以 $8^{-k/2}$ 衰减，轨道有限长。
\end{example}

\begin{example}[点收敛而无限长]\label{ex:dd2-half}
取上一旋转族的端点 $m=2$，即 $F(y)=\norm y\,Qy$。
仍有 $\Phi=0$、$\psi(\rho)=\sqrt\rho$，DD2 第一支取等，且
\[
 t^2=a^2+a^4,\qquad s=a^2,\qquad
 u_{k+1}-u_k=\frac1{1+t_{k+1}^2},\quad u_k=t_k^{-2}.
\]
非零初值的半径趋零，故 $x^k\to0$。但 $u_k\le u_0+k$，所以
\[
 s_k=t_{k+1}^2\ge\frac1{u_0+k+1},\qquad
 \sum_{k=0}^\infty s_k=\infty.
\]
这说明 DD2 的距离下降与点收敛本身不保证有限长度。
\end{example}

\subsection{比较级数与耦合长度势}

一般的距离递推 $t_{k+1}\le\varphi(t_k)$ 与步长界 $s_k\le B_\Phi(t_k)$
可通过比较级数求和。以下保留存在与留域条件：
$U$ 为开集，$U_R=\{x\in U:d(x)\le R\}$，$\mathcal J(x)$ 在 $U_R$ 上非空，
且其每个输出满足 DD2。约定 $\dist(x,\varnothing)=+\infty$。

\begin{theorem}[比较级数条件]\label{thm:dd2-series}
设 $\Phi(t)\le t^2/4$。取连续非减 $Q:[0,R]\to[0,R]$，
满足 $Q(0)=0$、$\varphi(t)\le Q(t)<t$（$t>0$）。令
\begin{equation}\label{eq:dd2-series}
 \overline B(t)=\sup_{0\le u\le t}B_\Phi(u),\qquad
 \Theta(t)=\sum_{j=0}^\infty\overline B(Q^{\circ j}(t)).
\end{equation}
若 $\overline B$ 有限、$\overline B(t)\to0$，且 $\Theta(R)<\infty$，
则满足 $\Theta(d(x^0))<\dist(x^0,U^c)$ 的 $x^0\in U_R$
所产生的每条允许轨道均可无限续接、留在 $U_R$ 并收敛到 $S\cap U$ 中一点，且
\begin{equation}\label{eq:dd2-series-tail}
 d(x^k)\le Q^{\circ k}(d(x^0)),\qquad
 \sum_{j=k}^\infty\norm{x^{j+1}-x^j}\le\Theta(d(x^k)).
\end{equation}
\end{theorem}

\begin{proof}
$Q^{\circ j}(t)\to0$，且 $\Theta$ 非减、
$\Theta(t)=\overline B(t)+\Theta(Q(t))$。
每步 $a\le Q(t)$，故 $s\le\Theta(t)-\Theta(a)$。
沿任意有限前缀求和，长度不超过 $\Theta(d(x^0))$；
严格初值条件使最后输出仍在 $U$，距离下降使其仍在 $U_R$，非空性保证续接。
再求尾和得式~\eqref{eq:dd2-series-tail}。有限长度给出柯西性，
距离趋零与 $S$ 闭性识别极限。以 $\Theta(R)$ 的可和项为控制，
还得 $\Theta(t)\to0$。
\end{proof}

该级数分别估计距离与步长，可能过粗。保留同一步中两者的联系，可以直接构造长度势。
设 $\psi$ 在 $[0,\rho_*]\subseteq D_\psi$ 上连续严格递增、$\psi(0)=0$，
区间外可任意。设 $R/2\le\psi(\rho_*)$，记
\begin{equation}\label{eq:dd2-q}
 b(u)=\lambda\psi^{-1}(u),\qquad q(t)=b(t/2),\qquad
 P(t)=\int_0^t\frac{u}{q(u)}\,du.
\end{equation}

\begin{theorem}[耦合长度条件]\label{thm:dd2-length}
设 $\mathcal J$ 在 $U_R$ 上非空，每个输出满足~\eqref{eq:dd2-assumptions}。
取 $D\ge0$、$D(0)=0$、$\Phi\ge-D$，并假设
\begin{align}
 &D(t)\le q(t)^2/4,\qquad 0<t\le R, \tag{F1}\label{eq:dd2-F1}\\
 &b(u)=o(u)\quad(u\downarrow0), \tag{F2}\label{eq:dd2-F2}\\
 &t/q(t)\le\overline g(t),\quad \overline g\text{ 非增},\quad
   \int_0^R\overline g(t)\,dt<\infty. \tag{F3}\label{eq:dd2-F3}
\end{align}
缩小 $R$ 使 $q(t)\le t/2$。若 $x^0\in U_R$ 且
$4P(d(x^0))<\dist(x^0,U^c)$，则每条允许轨道均可无限续接、留在 $U_R$，
并以有限长度收敛到某个 $x^\infty\in S\cap U$，且
\begin{equation}\label{eq:dd2-potential-tail}
 \norm{x^k-x^\infty}\le\sum_{j=k}^\infty\norm{x^{j+1}-x^j}
 \le4P(d(x^k)).
\end{equation}
\end{theorem}

\begin{proof}
若 $a\ge v>0$、$v\le R/2$，则 $s\ge b(v)$：
否则 $\rho\le s/\lambda<\psi^{-1}(v)\le\rho_*$，
原误差界与区间内严格单调性给出 $a\le\psi(\rho)<v$。
若 $s<q(t)$，取 $v=t/2$ 得 $a<t/2$，但
$t\le a+s$ 和 $q(t)\le t/2$ 又给 $a>t/2$，矛盾。故 $s\ge q(t)$，
由平方距离恒等式与 F1 得
\[
 t^2-a^2\ge s^2-2D(t)\ge s^2/2\ge q(t)^2/2>0.
\]
$q$ 递增，因而
\[
 P(t)-P(a)\ge\frac{t^2-a^2}{2q(t)}
 \ge\frac{s^2}{4q(t)}\ge s/4.
\]
零输入由 $\Phi(0)=0$ 驻定。沿有限前缀求和，由初值条件先得留域，
再用非空性续接，与比较级数定理的归纳相同。
若距离有正极限 $\delta$，则每步平方距离至少下降 $q(\delta)^2/2$，矛盾。
F3 保证 $P$ 有限且在零连续，望远镜求和给出尾界；
柯西性及 $S$ 闭性给出所述极限。
\end{proof}

取 $\psi(\rho)=\rho^{1/m}$，则 $q(t)=(t/2)^m$，
$t/q(t)$ 在零附近可积当且仅当 $m<2$。
因此例~\ref{ex:dd2-rotation} 在足够小邻域满足 F1--F3，轨道有限长。
另一方面，在比较级数中取本例的精确半径函数 $Q=\varphi$，则
$Q^{\circ k}(t)\asymp k^{-1/(2m-2)}$，
而 $\Phi=0$ 给 $B_\Phi(t)=t$；
$m\ge3/2$ 时式~\eqref{eq:dd2-series}发散，耦合长度条件仍适用。
这里省去的是对该比较级数的额外检验，F3 的积分条件仍然保留。
一般缺陷及正则性函数下实际有限长度的完整边界尚待进一步研究。
